\chapter{十九世纪：内乱与外来冲击}
\label{ch:nineteenth-century-conflicts}
\begin{figure}[htbp]
  \centering
  \begin{tikzpicture}[
    x=1cm,y=1cm,
    date/.style={font=\scriptsize\bfseries, text=dynasty},
    event/.style={draw=timeline, fill=paper, rounded corners=2pt, align=center,
      inner sep=3pt, font=\scriptsize, text=ink},
    note/.style={font=\scriptsize\color{source}, align=center}
  ]
    \draw[timeline, line width=1.2pt] (0,0) -- (12,0);
    \foreach \x/\year in {0/1851,1.8/1853,3.7/1856,5.4/1860,7.0/1861,8.7/1864,10.0/1868,12/1874}
      {\draw[timeline] (\x,.12) -- (\x,-.12); \node[date, below] at (\x,-.18) {\year};}
    \node[event, above] at (0,.42) {金田起事；\\永安建号};
    \node[event, below] at (1.8,-.62) {南京／天京；\\跨省战场展开};
    \node[event, above] at (3.7,.42) {天京事变与\\第二次鸦片战争};
    \node[event, below] at (5.4,-.62) {北京战事与\\新条约安排};
    \node[event, above] at (7.0,.42) {辛酉政变；\\总理衙门设置};
    \node[event, below] at (8.7,-.62) {天京陷落；\\战后重建起点};
    \node[event, above] at (10.0,.42) {捻军主力\\战争收束};
    \node[event, below] at (12,-.62) {日本出兵台湾；\\海疆问题再现};
    \node[note] at (6,-1.35)
      {战争中心的覆亡、条约程序、地方军事化和海疆危机各有不同的起点与余波。};
  \end{tikzpicture}
  \caption{十九世纪中叶冲突的编年定位（1851--1874，简表）。本图将太平、捻军、西北和西南战争的时段，同第二次鸦片战争、外交机构与战后军政重组并置，但不把它们归结为同一场战争。日期主要据《清文宗实录》《清穆宗实录》、军机处档、战事方略和中外条约；伤亡、军费、地方控制与恢复程度仍须依各地档案和多方材料分别核验。}
  \label{fig:nineteenth-century-conflicts-timeline}
\end{figure}

\begin{figure}[htbp]
  \centering
  \begin{tikzpicture}[
    x=1cm,y=1cm,
    place/.style={draw=timeline, fill=paper, rounded corners=2pt, align=center,
      inner sep=3pt, font=\small},
    court/.style={place, draw=dynasty, fill=dynasty!13},
    port/.style={place, draw=timeline, fill=timeline!13},
    war/.style={place, draw=source, fill=source!10},
    note/.style={align=center, font=\scriptsize\color{source}}
  ]
    \fill[paper] (-.65,-3.9) rectangle (7.85,3.85);
    \draw[dynasty, line width=.8pt] (-.65,3.85) -- (7.85,3.85);
    \node[font=\bfseries\large, text=dynasty] at (3.6,3.42)
      {1856--1864：条约口岸、京津战事与省际战争};

    \node[court] (beijing) at (3.75,2.55) {北京\\朝廷、交涉与\\1860年危机};
    \node[war] (tianjin) at (5.7,2.1) {大沽、天津\\1858--1860\\战事与换约};
    \node[port] (guangzhou) at (.1,1.45) {广州\\1856--1857\\战事、通商节点};
    \node[port] (shanghai) at (2.15,.45) {上海\\租界、军械、\\海关与商路};
    \node[war] (nanjing) at (3.9,.15) {南京／天京\\1853--1864\\战争中心};
    \node[war] (hunan) at (1.15,-1.45) {湘、赣、皖\\兵源、粮道与\\太平、捻军战区};
    \node[war] (northchina) at (5.7,-1.2) {华北、山东、河南\\捻军、河防与\\1855年河道变化};
    \node[port] (fuzhou) at (6.55,-2.65) {福州及东南港口\\船政、航运与\\条约网络};
    \node[war] (northwest) at (-.05,-2.8) {陕甘、西北\\1860年代战争\\与军需路线};

    \draw[->, source, line width=1.1pt] (tianjin) -- node[above,sloped,note]
      {进逼北京的军行} (beijing);
    \draw[->, timeline, dashed] (guangzhou) to[out=-15,in=155] node[above,note]
      {沿海交涉、航运与消息链} (shanghai);
    \draw[->, timeline, dashed] (shanghai) to[out=0,in=-100] node[right,note]
      {江海交通与战时物资} (nanjing);
    \draw[->, source, line width=1.1pt] (nanjing) -- node[left,sloped,note]
      {长江中上游战事、筹饷与征发} (hunan);
    \draw[->, source, dashed] (nanjing) -- node[right,sloped,note]
      {华北战场与河防的不同节奏} (northchina);
    \draw[timeline, dashed] (shanghai) -- node[right,note]
      {港口间航运联系（示意）} (fuzhou);
    \draw[source, dashed] (hunan) -- node[below,sloped,note]
      {西北军务与地方转运} (northwest);

    \node[draw=source, fill=paper, rounded corners=3pt, align=center,
      text=source, font=\small] at (6.15,.65)
      {铜色节点：口岸／交涉空间\\灰褐节点：战争、河防或军需空间\\两者都不是连续控制区域};
  \end{tikzpicture}
  \caption{第二次鸦片战争、条约口岸与内战的相邻空间（1856--1864为主，示意）。图将广州、上海、福州等口岸及京津交涉节点，同南京、湘赣皖、华北和西北的战争、河防与军需节点并列，说明条约签署、港口制度、沿海航运和各省战事具有交错但不同步的地理节奏。依据《清文宗实录》《清穆宗实录》、军机处档与战事方略，《筹办夷务始末》、中外条约及照会文本、总理衙门和海关档；并参照地方志、督抚军需档、上海租界记录、英法外交军政档与港口报刊。图不列举完整口岸清单，路线不按比例；条约文本的签署、批准、换约和地方执行须分别核验，港口、租界、战区或箭头均不表示周边地区的统一控制、人口态度或贸易覆盖。}
  \label{fig:treaty-ports-wars}
\end{figure}


\noindent 两幅图把一段常被分开讲述的历史放在同一页上：一边是广州、天津、北京和条约口岸的海路战事与交涉，另一边是南京、湘赣皖、华北及西北的征发、河防和军需。它们相邻，并不意味同一命令、同一条航线或同一笔款项能同时抵达这些地方。图中的线条只提示可被争夺的联系；战区、口岸和条约文字都不能替代对一处城镇、一段道路或一个社群实际处境的判断。十九世纪中叶的危机，正是在这种不同速度的联系中展开。

\section{1851—1856：从金田到天京，战争改变了征发的尺度}

咸丰元年，\eventanchor{EQ-1851-JINTIAN-UPRISING}{金田起事（1851）}使广西的团练、交通和地方军政迅速卷入一场将延及多省的战争。太平军北上经湖南，1853年夺取南京并改称天京，随后北伐、西征与长江沿线的争夺同时展开。城池易手不只是旗帜更换：粮道、船只、厘金、逃难路线和地方武装的选择，都决定哪一支军队能在某段河流或某座府城继续停留。

清廷不能只依靠原有绿营和八旗解决这场战争。湖南团练和湘军雏形的组织，以及此后各地勇营、捐输与饷源网络的生长，使地方官绅、将领、商人和运输中介获得更大的筹饷、用人和调兵空间。这并非中央命令忽然失效，而是命令若要成为一支能行军的队伍，仍须穿过省内外的银粮、船运、担夫和信用关系。一笔已经获准的饷银，要折色、收纳、解送、交收；一船米粮则还要有可通的江面、肯装运的人和不被两军扣留的渡口。这样的接力提高了若干战区的动员能力，也把军费、任命与武装关系更牢地嵌进省级网络；对沿线家庭而言，征购、雇役、避难和借贷却往往没有这样清楚的制度边界。

1855年铜瓦厢决口改道与捻军活动扩大，使华北、淮北成为同长江战场相连却节奏不同的危机地带。河工请款可以显示朝廷和地方要修防什么，不能证明堤防立刻奏效；军报可以显示兵力被调往何处，也不能替一村一镇说明粮价、迁徙或征发怎样落下。同年，\eventanchor{EQ-1855-PANTHAY-OUTBREAK}{滇西多地冲突扩大}，矿业、商路、官军冲突与社区关系的紧张又在西南形成另一条战争链。云南军报将许多关系置入征剿和行政分类的语言，地方志与碑刻留下的又是断续的地点和人名；它们足以提醒我们大理及其周边并非静止的边角，却不足以把不同社群写成一个意志一致的阵营。把这些地方一概称作内地战局的回声，会抹去各自的道路、市场和行动者；把它们只写成孤立的地方乱子，也会看不见全国军力被牵制后留下的空隙。1856年天京事变削弱太平天国的内部协作，却未让战争的社会代价或地方选择变得单线可读。

\section{1856—1864：两场战争在不同道路上相遇}

第二次鸦片战争在1856年爆发，英法军队的海路行动、广州失守、天津谈判和大沽口战事与内战并行。1858年\eventanchor{EQ-1858-TIANJIN-TREATIES}{《天津条约》等的签订}扩大了通商、使节与传教的制度语言；它不是沿海、内地和边地生活在同一天被改写的结果。签署、批准、换约、翻译、地方知会和实际执行属于不同环节，且须由不同的文书和地点来检验。一份中外对照条约能够让读者追到约文如何措辞，却未必保留译者、地方官和港口商人怎样理解或回避它。口岸的航运和租界网络能为军械、消息和人员流动提供条件，却不能被写成一条从海岸直达每个战场的输送带。

1860年\eventref{EQ-1860-BEIJING-CAMPAIGN}{英法联军进逼北京}，圆明园遭焚掠，咸丰帝出走热河；其后的\eventanchor{EQ-1860-CONVENTION-OF-PEKING}{《北京条约》及相关中俄条约}把外交、边界和通商条件固定在文本中。文本可以确定国家作出何种承诺，不能单独说明一座港口、一段边界或一处教案怎样被执行，更不能把东北的条约文字直接当作当地控制或居民认同的地图。京师危机同长江的围城、华北的河防并无自动的单一因果，却同时挤占了朝廷的船路、饷项、官员注意力和递送信息的能力。

1861年咸丰帝去世，辛酉政变重组了宫廷权力，\eventanchor{EQ-1861-ZONGLI-YAMEN}{总理各国事务衙门的设立}随即提供了处理持续交涉的常设机构。它是危机后的制度应答，不是一个已经能替地方完成外交、海防和财政事务的全能机关：经费从哪里筹来，文本由谁翻译，命令经由哪些督抚、口岸与中介传递，仍决定其作用的边界。与此同时，厘金、捐输、借款与各省协饷把不同前线的开支接到商路和地方税源上；“筹到”与“送到”之间仍隔着折算、押运和验收。1862年以后上海的军火、租界网络和常胜军进入江南战场；这些资源改变了某些战场的条件，却不应被写成外来技术自动决定胜败的故事。它们能够抵达何处、由谁付费、在何种地方关系中被使用，同样是战争的一部分。

1864年湘军攻陷天京，\eventanchor{EQ-1864-TIANJING-FALL}{太平天国核心政权覆亡（1864）}。这是一个明确的军事终点，不是人口、土地与行政关系立即复原的起点。围城、逃亡、征发和失去生计的经历不可能由一份战报或一项善后章程加总；城中与乡村留下多少人、何时能耕作或重开交换，须回到保存不均的地方材料逐处追问。田契、账簿或地方志偶尔能照见归户、典卖和市场重启，却更容易留下能书写、能保存文书的人；不能把这些断面当作整个江南的复原曲线。捻军仍在华北和淮北活动，陕西、甘肃及其后的西北战争又把军队、商路、不同社群和地方社会拖入新的多年冲突。战后设局、厘金和军工项目确实扩大了国家及地方官员的行动工具，也把财政和军权更深地留在省级网络之中：它们解决了眼前调度的难题，却把成本以不同方式留给了驻军地、转运线和纳税、借贷或服役的人。

\section{1865—1874：重建不是回到战前}

1865年以后，重建的地图很快越出江南。陕西、甘肃的长期战争与喀什噶尔政局重组，使河西走廊、绿洲城镇、山口和中亚商路成为军队与地方社会都必须重新计算的空间。清廷后来向西推进时，面对的并不只是一个可在地图上收回的空白，而是已被守城、征粮、水利、商队与跨境观察重新组织的关系。西北官书的分类不能概括这些地区的人群及其选择；清方军报、俄国观察与维吾尔文、察合台文地方文书所见的范围和政治目的也并不相同。它们常连地名、人名和纪年也未必直接相合：一份军报所称“克复”的城镇，不能不经路线、驻守和地方材料的核对便延展为一片稳定控制。正因如此，饷运的记录能让我们看见筹划和调度，却不能把到站、交接、征购或地方承受力一概视作既成结果。

福州船政局、江南制造局等项目在1860年代出现于战争与口岸交错的环境。它们需要经费、译书、技术人员和地方合作，因此不宜以“洋务”一个词代替差异甚大的实际进度。厂局设立、预算呈报或一项技术学习计划，首先证明的是政策和组织意图；产能、军用效果以及对沿江、沿海社会的影响，还要依账册、具体船只和地方记录分开考察。1860年代末捻军战事的收束与西北战争的扩大并行，正说明所谓战后财政并非从军费中解放出来，而是在不同前线、港口与省份之间继续选择轻重。

1870年天津教案使教会活动、地方传闻、法国领事保护和清廷交涉在一座城市相撞。照会和处分文书能够保存各方怎样提出责任，不能替代无从留下文字的居民经验，也不能把城市中的教民、传教士、官员和民众分别写成意志一致的阵营。1871年牡丹社一带的事件及其后\eventanchor{EQ-1874-JAPANESE-TAIWAN-EXPEDITION}{日本出兵台湾}，又把琉球船民、台湾南部原住民地区、地方防务和对日交涉牵进同一轮紧迫的估算。日方以事件提出的政策主张、清廷能够宣称的管辖范围、以及当地社会实际可被动员或保护的程度，不能相互替代；海峡与海岸由此不再只是条约口岸之外的边缘。


\section{战争的续航：银、粮、船与人的接力}

若从1851年金田起事一路看到1864年天京陷落，最容易被连续的战线遮住的，是战争每天都要重新取得的物资条件。军队并不因为奉有上谕、挂出名号或占得城池便自然拥有粮食。米要从产地或市镇收买、征取、捐输或借得，装入仓廒，换成船运或陆运，经过验收、转交和损耗，才可能在营前成为一顿饭；银则在平码、折色、汇兑、解库和分发之间改变形态。账簿中的一笔协饷既不是一袋已经抵达的银子，也不是某个兵丁必已领到的月饷。它首先说明某个衙门、捐户或商号被置入一条筹措和交割的关系；在道路被截、平码变动、胥役拖延或地方另有急用时，原定的下一环完全可能断开。把这种断开视为例外，反而会误解这场长年战争何以既能延续又常常陷入欠饷、哗变、逃散和临时征购。

湘军及其后各省勇营的扩张，使筹兵与筹饷越来越紧密地系于省级和跨省的人际网络。督抚、将领、州县、绅士、盐商、平码局、牙行和船户并非同一种行动者，也没有共同的利益。有人愿意垫款，是为了保全商业、取得声望、维系官场信用，或仅仅因为已被要求承担；有人接运，是在兵役、风险和报酬之间作出有限选择；有人则以逃避、拖延、转移存粮或诉诸亲族来抵抗负担。所谓捐输在档案中常有成数和名目，却不应因此把地方社会写成自动响应的财源。捐例、保甲、团练和厘金所能触及的范围，要看地方文书、契约、账目和个案中是否留下实际交纳、拖欠、抵押或争讼的痕迹。保存下来的材料又偏向能同衙门或市场交涉的人，贫户、流民和被强迫服役者的经验更难被完整记录。

长江是军事路线，也是被战争反复改写的市场空间。南京成为天京以后，上游粮道、湖口、九江、安庆和皖南、赣北的水陆节点都不再只是地图上的地名。哪一段江面可通，哪一个渡口能征得船，哪处仓房可暂存米盐，取决于季节、水位、兵舰、地方武装和沿岸居民的应对。大宗运输的叙述很容易把船夫化为无声的工具；实际上，船户要面对被征发、被扣留、遇劫、涨价和家计中断的风险。清军、太平军及地方团练对船只和粮食的索取，也不必然按同一规则进行。军报中“筹办”或“解运”的动词，是行政目标的语言；要判断它在何处形成实际运输，仍须考察平码、船数、到站、交收和沿途损耗等彼此独立的记录。

战争还重排了信用。银根紧张或实银难以移动时，票号、钱庄、商帮和熟人之间的担保能使一部分款项跨越省界；但票号的汇兑网络并不是无摩擦的国家管道。它依赖分号、信使、信誉、平码和地方秩序，而这些条件恰是战争所侵蚀的对象。军需借款可以显示官军进入了商业信用的范围，不能据此推称财政已经现代化，或说商人能够自由决定国家军事。相反，垫款者可能因征调而被卷入，承办者也可能通过信息和手续获得新的位置。县城里的典当铺、米行和银炉所感到的，不一定是京师所称的“筹饷”，而可能是货币难换、谷价波动、熟客失踪和债务链条加长。

1855年黄河铜瓦厢决口之后，华北的水系、堤防和道路问题又给这种接力增加了沉重的一层。决口与改道不是一条单向的自然背景，而是使田地、村落、渡口、盐运、漕运遗绪和军事移动重新发生冲突的过程。河工奏报里有估工、请款、挑筑和堵口的程序，地方志和赈务档案或能保留若干受灾地点；但同一份“工竣”呈报不能证明每段堤防都稳固，更不能替代沿河家庭对失地、借贷和迁徙的叙述。捻军的活动在这样的地带扩大，也不能简化为洪水或饥困必然产生的军事结果。水患、赋役、地权、武装流动和地方结社相互作用，却须在具体县份和年份中辨认其先后，不能把宏观相关写成单线原因。

华北和淮北的战场说明，兵员并非只从军籍上产生。保甲、团练、雇勇、逃兵、被俘者和随军家属在不同情形下进入或离开武装队伍。一个人被县衙列作丁壮，可能正同时是佃户、挑夫、盐贩、父亲或逃难者；名册中的“可用”也未必等于能够离乡。地方官为筹集夫役而开出的名额，能够证明他所期待的供给，不能证明每一户都按额出人。途中死亡、逃散、替代、贿免和雇募的细节常被汇总数字抹平。对地方社会来说，战争的压力不只在战场的杀伤，还在劳力被抽走之后的耕作、照料、婚姻与债务如何无法照常安排。

省际协饷常被当作晚清权力下移的简洁证据，实际却更像一组不断重谈的关系。中央仍能以谕旨、部议、任命和考成塑造框架，地方督抚也须在朝廷的授权、同僚的协商、商人的信用和本省民力之间调度。某省“应解”的数字，既不是中央能随意调用的现银，也不是地方已经完全自主的收入；它是一个可能被折抵、缓解、催提、挪用或重新分派的行政和财务要求。厘金的设置同样如此。关卡、平码、查验和征收人员为战费提供了可能的来源，也增加了货物过境的时间和成本。商路没有因为征厘而停止，地方也没有因此成为同质的税收机器；不同货类、不同口岸、不同省界的规避、谈价和冲突，才构成税制实际的形状。

沿海口岸的关税、航运和外国商行则使战争经济多了一条与内陆并行的通道。上海、宁波、福州、汉口等地的贸易信息、保险、轮船和买办网络能让货物、金属、军械和消息更快移动，却并不把沿江腹地变成可以任意调用的后方。口岸报关数字说明在特定制度和口径下登记的货物价值，不能直接换算为军队获得的物资；外商的合同、领事的照会和租界的记录也首先说明其参与者所关心的风险和权利。1860年代上海的军事采购与常胜军活动，显示部分外来资源在江南战局中具有可见影响，但采购、付款、翻译、训练、弹药保管和战场指挥之间仍有层层条件。把这些资源写成某一方胜败的唯一钥匙，会同时抹去地方军队、居民与市场承担的成本。

\section{离乡、归乡与留在原地的人}

战争年代的迁徙不能只用“逃难”二字概括。有人在城池易手前随亲族先走，有人在征粮、疫病、洪水或田租压力下分批离开，有人沿着亲戚、同乡、会馆和商路的线索前往城市、口岸或较少受战事波及的县份，也有人在渡口、寺庙、义仓和亲属家中短暂停留。还有许多人并未离开，却在田地荒芜、市场中断和劳力减少的处境中被迫改变作物、借贷方式和家庭分工。地方志中“流民”“难民”这样的总称，帮助我们看到危机规模，却会掩盖迁移距离、性别、年龄、财产和返回机会的不同。能留下呈报、契约或日记的人，往往比无籍、无资或被迫流动者更容易进入史料。

江南的城镇与乡村并非分别经历战争。城内需要粮柴、劳力、船运和信息，乡村则要应对过境军队、征购、躲藏和市场需求；二者通过河港、墟市、宗族、雇工和借贷紧密相连。围城会改变周边地区的粮价和道路选择，城市失守又可能使人口、货物和政治传闻向邻县移动。但不能从一个大城市的遭遇推演整个区域。苏南、皖南、赣北、湖北和湖南的战事节奏、作物结构与水路条件都不同，即使同一府内，靠近大路的村庄和偏远山地所面对的征发也可能相去甚远。把“江南毁坏”当作解释一切的背景，容易将具体家庭的韧性、损失和机会一并抹掉。

土地关系在战后尤其需要避免用复垦与买卖的几个例子讲成整齐的复原。荒田被谁耕种，原业主何时返回，佃户能否续租，典当债务如何清理，宗族祠产和庙产是否仍能维持，这些问题往往在不同村镇有不同答案。田契能保留交易的双方和价额，却未必写出被迫出售的原因；赋役册可显示重新纳入征管的企图，却不说明每户的实收；地方志喜欢以“招徕”“垦复”描述行政成效，也可能压缩了争地、占佃和流民定居的摩擦。重建不是恢复一张战前的土地清册，而是让新的占有、租佃、担保和亲属关系在不平等的条件下重新稳定或继续争执。

妇女、儿童和老人更常以家属、烈属、寡妇、孤儿或被赎者的身份零星出现于档案，正因如此不能把他们放在战争叙事的边缘。男性劳力被征调、死亡、逃散或长期随军时，留守者要承担耕作、抚养、借贷和与衙门、宗族交涉的责任；但现存材料多由官员、士绅或男性亲属书写，难以直接听见其判断。贞节旌表、赈济名册、诉讼和慈善账册有时能照出家庭破裂后的制度回应，同时也以自身的道德和行政分类塑造了可被看见的人。我们可以据此说明战争重压落入家庭内部，却不应把沉默解释为顺从，也不能从少数显名案件推广为所有妇女的共同经历。

宗族、庙宇、会馆和善堂并非天然的救济机器。在一些地方，它们可能筹米、安置流民、修桥、收埋遗骸或调停纠纷；在另一些地方，财产损失、成员离散和派系冲突使其无力行动，甚至成为争夺资源的场所。善堂账簿记录了施给者所愿意登记的对象，不等于实际所有受助者；碑记的捐资名单体现名望与纪念，也不必然等于工程已经完工或救济长久有效。正是这种有限性使它们珍贵：它们让我们看到国家命令之外，地方社会如何试图把粮、钱、棺木、药品和劳力组织起来，而不是把“民间”抽象成一个同样整齐的主体。

战争中的宗教空间也不能只以教义或政治阵营说明。寺观、教堂、清真寺、祠堂和会馆既可能是避难、储粮、集会和传递消息的地点，也可能因其财产、象征或人员关系而卷入冲突。官方文书对“会匪”“教案”或“回”的称呼服务于治理和追剿分类；外来传教记录则可能以受害、保护或开化的语言组织经验。两种文本都不足以把当地信众、居民和武装者写成同一意志。需要追问的不是抽象的“宗教导致冲突”，而是在何时何地，土地、司法、贸易、保护关系和传闻如何使某处空间成为可争夺的节点。

人口数字尤其要求克制。战乱后的户口减少可能反映死亡、逃亡、漏报、并户、行政停顿或丈量口径变化；地方志的灾损数字、传教士的估计和后出研究的推算各自有覆盖范围。它们可以共同提示人口流失和恢复并非小事，却不能在没有区域、年份和算法说明的情况下合并为一个精确的全国总数。更稳妥的写法是把数字放回生成它们的机关和目的中：赈务奏报为何需要估计灾户，赋役清册何时恢复，城市慈善组织向谁登记，海外观察者从何处得到消息。数字既是受灾经验的痕迹，也是行政与舆论争取资源的工具。

1860年代后，有些地区出现归乡、招垦、移民和新集镇的迹象。这些变化并不必然意味着旧秩序原样复归。新来的佃户、驻军家属、商人和手工业者会带来新的租约、市场和权利要求；原有居民则可能借亲族、地缘或文书证明主张。道路修复、桥梁重建和渡口恢复也使货物重新流动，却可能改变哪一处市镇成为中心。地方重建的时间因此不同步：一座县城修复城墙，与周围村庄恢复耕种、与某一商路重开、与流散者能否获得安身之地，是彼此关联而不能互相代替的过程。

\section{西北、西南与边路上的长期战争}

1860年代的西北战争把陕甘、宁夏、青海边缘、河西走廊和新疆的许多地方卷入长时间的军事与社会重组。使用“回乱”“回民军”或“平乱”等后出的概括时，必须记得这些标签来自不同政权、行政和叙事传统，不能代替当地社会的自称、联盟和分歧。村庄、城镇、商队和宗教网络中的人们并不因被同一类目归档便作出相同选择；清军、地方武装、商人、乡绅及各社群内部的领袖和普通人也都处在彼此变动的关系中。冲突的暴力、报复、迁徙和饥困是真实而深重的，但若把它们写成恒久的群体仇恨，就会掩盖行政失序、军粮、地权、债务、地方防卫和跨区域消息的具体作用。

西北军需的困难在距离上格外可见，却不能用“遥远”一句带过。由关中出发的银粮要经过道路、驿站、山口、河谷和市场，随季节、天气、牲畜和安全状况变化。驮运所需的骡马、骆驼、草料和赶脚人本身也是要筹措的资源；一支部队向前推进，意味着后方不断有人在修路、运粮、看守仓栈和照料牲畜。奏折中列出的运道和运价能够显示官员如何估算这些问题，不能证明粮食按数到达，亦不能把沿线居民的出售或被征购理解为自愿。军需压力会把市场和家庭推入战场，但市场参与者也会依据风险、价格和关系作出自己的取舍。

陕西和甘肃的重建因此不能只以“清理战场”概括。荒地、废城、寺产、商铺和水利设施的修复涉及谁有权回归、谁可领垦、谁承担赋役以及谁能得到官府承认。移民政策和招垦文书记录的是国家希望恢复人丁、田赋和秩序的方案；它们不能自动证明移民规模、土地质量或新旧居民之间的相处。地方志、契约、墓碑和口述线索有时能补足某一处社区的变化，但同样受保存和后来的叙事选择限制。对于没有留下文书的人，我们只能承认材料的缺口，而不是用“安定”或“恢复”填补沉默。

云南自1850年代中期扩大并延续至1870年代初的冲突，则把矿业、马帮、山地交通、城镇商业和不同社群关系拉入另一种节奏。滇西的铜、银、盐与跨境贸易并非静态的经济背景；矿山的劳工、税收、债务和商路保护本就容易成为地方权力的焦点。大理及周边政权和军队的形成、对抗与消长，不能仅从昆明或北京的奏报来理解。清方军报、地方文献、商人记录以及来自缅甸、中南半岛方向的观察，各自所见的路线与人群不同。它们并读可以防止把云南写成中原战场的附庸，却仍无法让我们精确量出每一处村寨的损失或忠诚。

滇西的军事推进同样依赖粮道，而粮道又会改变地方社会。山地道路可否通行，取决于雨季、桥梁、脚夫、驮畜和沿线市场；从一份“克城”军报到稳定驻守之间，仍隔着守军的给养、疾病、当地协作和后续行政。官书以收复、招抚、剿办排列事件，是一种从政权中心向外看的时间；对于被迫离散的人，时间可能是播种季节错过、亲属未归、债券到期或道路重新可走。将两种时间并置，不是削弱军事史，而是避免胜败叙事把社会后果当作已经解决的尾声。

新疆与中亚的情势更提醒人们，边界不是地图上一笔完成的线。1860年代喀什噶尔及南疆各城的政局变化、阿古柏势力的建立和扩展、俄国在伊犁的行动，分别牵动绿洲城镇、伯克与宗教人士、商队、农民、军人和跨境观察者。汉文官书、俄文外交材料、英印观察以及维吾尔文、察合台文或波斯文材料所用的地名、纪年和政治词汇并不完全相同。把它们简单译为同一套国家主权术语，会造成一种并不存在的透明性。尤其是地方文书的存留不全、版本和释读有争议时，任何关于民众态度或经济规模的断言都应保留出处和范围。

左宗棠主持的西征与1870年代新疆战事，常被叙述为财政与后勤的胜利。这种说法若不拆开，便会把筹议当作实现。军饷来源、海关收入、借款、协饷和粮运确实构成可追索的政策与经办网络，但每一段都有等待、折损、拖欠和地方负担。军队推进到某城，不能直接证明周围绿洲、山口和牧地已经纳入稳定治理；册封、设官和驻防也只是行政行动的起点。要讨论控制的范围，应分别询问驻军是否持续、税役如何征收、水利由谁管理、商路如何通过，以及地方文书和外来观察是否留下相互约束的证据。没有这些区分，“收复”就会被误当成一个瞬间完成的社会事实。

1871年俄军进入伊犁，使西北问题同时具有军事、外交和地方生计的层次。俄方的占领记录、清廷的交涉文书与当地社会的材料都不会自然说同一种话。外交照会能确定一方提出的主张和谈判的时序，却不能代替伊犁居民在税收、贸易、迁徙和安全上的经验；地方传闻和回忆能提供另一种时间感，也须辨别写作时距和传播路径。边地的“归属”在条约、驻兵、贸易、田地和人身流动中被不同方式实践，不能把后来的结局倒投为当时每一个人的确定认知。

1870年代中期华北的旱灾与饥荒，又使北方的军事和财政压力同救荒、运输和迁徙相交。灾情奏报、义赈账册、报刊通信和外国救济记录各自能保存粮价、募捐、运输和受灾地点的片段，但都不等于全面统计。赈款拨定、粮食起运和实际到达之间仍有道路、仓储、侵吞、转卖及地方分配等问题；一份感谢或结案文书也不证明没有遗漏。饥荒中的流动者未必能进入名册，留在村中者也未必得到同等救济。把灾荒看作与战争无关的自然插曲，会忽略此前征发、道路破坏和市场重组所留下的脆弱性；把它完全归为战争后果，同样会抹去气候、水利和区域粮食差异。

\section{条约、口岸与并不相同的外国行动者}

第二次鸦片战争的海上行动与内战并行，既不意味着外国军队能够任意支配清朝全部疆域，也不意味着内陆战争同海岸无关。英法军舰、商人、传教士、翻译、领事、买办、海关雇员和其他外国观察者有不同的机构、利益和信息来源。“外国人”不是一个能解释政策或行动的单一主体。英国政府的外交意图、法军将领的军事判断、商行的风险计算、传教网络的地方关系和报馆的新闻选择，常在同一城市中彼此不尽相同。中方的朝廷、总督、地方官、商人、乡绅、雇工和居民也一样。只有把行动者拆开，条约口岸才不会被写成外力与中国这两个整块相撞的舞台。

1858年天津条约及1860年北京条约的文字，提供了理解通商、使节、传教和赔款安排的重要起点。条约的中文、外文版本、换约程序和照会来往可以让我们追问何种权利被提出、如何表述、由谁承认；然而签字或批准并不自动等于某一县份、教堂、港口或内地道路已经按文运行。翻译用词、地方告示、领事保护、诉讼实践和民众传闻会让同一条款呈现不同面貌。对“治外法权”“传教权”或“通商权”的叙述若只停在条文层面，容易把地方争端的土地、债务、刑案和人际保护关系消失掉；若只叙地方冲突，也会看不见制度语言如何改变当事人能够诉诸的渠道。

广州在1857年失守、天津与大沽口的战事、1860年北京危机，都显示海防与外交的时间并不由同一节拍控制。军报记录进攻、退守和调兵，外交文书记录交涉、换约和赔款，城市材料则可能留下停市、避难、雇佣和谣言的线索。它们能相互校正某些日期，却不该被拼成一幅全知的图景。圆明园被焚掠的物质损失与象征冲击可以通过中外记录和遗存讨论，但个体参与、命令链和所见范围仍需区分证据层次。将京师的危机直接解释为所有地方政策即时转向，或将沿海行动直接解释为内战结局，都会越过文书和道路之间实际存在的距离。

1861年总理各国事务衙门的设立，显示清廷试图为持续的外交事务建立较稳定的处理机制。机构设立能够证明一个新的办事框架获得正式承认，不能证明它已经拥有统一的预算、翻译能力、信息网络或对督抚的直接指挥。照会到达京师、译成何种文字、交由谁拟答、再送回口岸，需要时间和人员；地方督抚在海防、通商和教案中的报告也会带着本地资源和风险的考量。把这个机构称作“近代外交的开始”可以作为制度史的概括，但若不说明它在何种问题上、通过何种文书发生作用，就容易把日后的组织能力提前投射到1860年代。

海关制度的变化和外国雇员的参与，同样不应只从“控制”或“利用”的二分理解。税则、报关、巡查、船期和统计使口岸贸易获得新的可见性，也创造了新的中介位置和利益冲突。海关统计的单位、年度和登记范围须被逐项说明；它记录经过特定关口的贸易，不能等同于全国生产，更不能把走私、内河转运和未登记交换化为零。海关收入被用于不同财政需求的安排，是可以在预算、奏折和账目中追查的事项；至于一笔收入何时、以何种折色到达何处，仍需有交接和支出的证据。

传教活动在条约之后的扩展也不能用一条直线来叙述。教士进入某地、购置房屋、设立学校或诊所，可能依赖条约语言、领事保护、地方官许可、华人教徒的关系和地方市场；每一个条件都可能受阻。教案中的传闻、儿童收养、土地交易、司法管辖和地方竞争，往往比抽象的宗教对立更能解释一场冲突何以在某时某地爆发。1870年天津教案的中外材料有助于追踪交涉、处分和舆论，却不能让我们把城市居民、教民、官员和外国人各自缩成一致的集体心理。新闻报道尤其要辨认首发、转载、译述与立场，不能把报上的即时判断当作已经核实的事实。

台湾南部在1871年牡丹社一带的事件及1874年日本出兵之后，显示海洋边界、琉球船民、原住民地区和清朝地方治理相互交错。日方的出兵理由、清廷在交涉中的表述、台湾地方官的防务安排和当地社群可被看见的行动，并非同一层次的材料。行政上称为“化外”的区域，不是没有社会、道路、土地关系和政治秩序的空白；相反，正因清朝的行政触及不均，任何单一的主权声明都不能替代现场关系。原住民自述保存稀少，更要求写作者避免把官府、日方或汉人移民的叙述直接当作其意愿。海图、条约和军舰航线可说明外部目光如何绘制海岸，却不能替代山地、社群和口述传统中的历史。

日本出兵后的交涉还使琉球问题的不同位置更为显露。琉球王国、萨摩及日本中央政策、清朝册封与海上贸易关系在十九世纪并非可以由一个后来国家边界完全概括的对象。对船难、赔偿、保护责任和管辖的争论，既是外交文本的问题，也牵涉海上航行的风险和地方接触。若只从北京或东京的照会理解，便容易把岛屿、港口和南部台湾的行动者排除；若只从一宗地方事件理解，又会看不见它如何被纳入更广的外交竞争。两种尺度都必须保留，且不能把1874年的结局预设成后来主权争论中唯一可能的含义。

\section{1875—1877：恢复的限度与叙事的边界}

到1875年至1877年，许多地方已经不再处于太平天国战争最激烈的阶段，但“战后”并不是一个所有地区共同进入的时刻。江南的复垦、华北的河工和赈务、云南的善后、西北的军运和新疆的战事同时存在。中央与地方的财政选择仍面对彼此竞争的需要：修堤、赈饥、给饷、偿债、办厂、设防和维持常规行政都需要钱粮与人力。将其中任何一项写成国家能力已经恢复的凭据，都会忽略其余项目可能被推迟、折减或转嫁。重建中的“恢复”常是一种比较性的判断：与最坏的年份相比，市场或行政或许重新运转；与战前相比，人口、地权、道路和债务却可能已是另一种结构。

福州船政局、江南制造局等项目在这个背景下更应被看作具体的组织实践，而不只是时代标签。厂局的选址、经费、人员、翻译、机器采购、试造和产品去向各有不同的时间。开局章程、预算呈报和官员奏议可以证明某种规划得到提出或批准；设备是否到场、工匠如何训练、船只是否按期完工、产品如何被使用，则要依赖账册、检验、航运和地方材料进一步核对。技术学习不等于技术能力立即转化为军事优势，外国雇员的存在也不等于他们控制全部决策。把制度意图、生产过程和战场效果分开，才不会以一份奏折替代多年累积的工作与摩擦。

1870年代的报刊提供了新的公共信息渠道，却不应让研究者误以为社会因此突然变得透明。《申报》以及其他城市报纸能保存灾情、募捐、市场和外交消息的日期化叙述，也能让我们追踪信息如何被转载、质疑或商业化；它们主要服务于特定城市、读者市场和通讯网络。电报和轮船缩短某些信息与人员的移动时间，但消息的发出、收取、翻译、编辑和再传播仍会产生延迟与变形。报刊中可见的“舆论”是发声者、编辑与读者之间形成的局部空间，不能被直接称作全国民意，更不能替代未能进入印刷媒介的乡村经验。

对这一时期的战争和重建，最可靠的写法不是寻找一个包容一切的解释，而是让问题同材料匹配。《清文宗实录》《清穆宗实录》与军机处档可帮助确定朝廷何时接到何种报告、发出何种命令；方略与军报能追索清方如何排列战事和奖惩；地方志、契约、账簿、碑刻和日记可以在有限地点上显示市场、土地、迁移和救济；条约、照会、领事材料与报刊能追踪跨国交涉和消息；多语文书、城址和物质遗存则为边地提供不完全但必要的另一种尺度。它们不是可以互相替换的证词。材料越丰富，越应说明其产生机构、语言、时间、保存状态和沉默的对象。

尤其需要警惕用官方动作替代社会结果。上谕证明命令被发布，奏折证明某人向上级陈述，条约证明缔约方在文本上同意，工竣册证明工程在某种程序中被宣布完成，战报证明报送者愿意呈现的战况。这些都具有历史价值，却不能自动推出地方服从、军饷到手、道路畅通、居民认同或秩序稳定。反过来，一份日记、传教报告或地方传闻虽能使宏大叙事出现裂缝，也不能独力代表整个地区。将命令、实施、抵抗、结果分别置于可核查的证据下，并承认某些环节无从确定，不是削弱叙事，而是让战争中的行动者不被事后结局吞没。

1851年至1877年的冲突由此不是从内乱通向自强、从危机通向恢复的一条单行道路。它包含多条相互交叉的道路：长江上的船运与围城，黄河改道后的堤防与流民，陕西甘肃的军粮与迁徙，云南的矿路与城镇，新疆绿洲与中亚交涉，天津和北京的条约程序，台湾南部的海上接触，以及各地家庭为活下去而作出的沉默选择。它们有时被同一份财政、同一条电报或同一轮外交牵动，却不因相连便在同一时刻发生同一种变化。保留这种不同步，才能既看见帝国在危机中重组的能力，也看见这种能力如何建立在被征发、被迫迁徙、承担风险和难以留下文字的人们身上。

\section{多重前线之间的行政时间}

把十九世纪中叶的战争只理解为中央与地方之间的权力转移，也会遗漏行政本身所经过的时间。北京发出的谕旨要由驿递、书吏、地方官和军营接续；地方递出的奏折则经过起草、转递、抄录、批答和再行文。紧急军情可以加快若干环节，却不能消除道路、天气、战事与人手造成的间隔。命令抵达时，当地的粮价、兵力、河水和人际关系可能已经改变。于是，同一道命令在不同省份并不是同一时刻、同一资源条件下被理解和处理。档案里相邻的朱批和奏折，常使这种间隔看似消失；把它重新放回通信和执行的路径中，才能理解为何中央的意图、督抚的陈述和县镇的处境会彼此不同。

这种行政时间也塑造了“战功”的可见性。将领的奏报需要把零散的侦察、交战、俘获和补给组织成可供上级判断的叙述，因而会突出可以奖惩、可以归责和可以纳入整体战局的事。方略与实录在后来的编纂中又会进一步排列这些报告，使若干节点成为明确的转折，而漫长的守备、疾病、逃散、谈判和地方自保则更容易缩成背景。并非因此应当否定军报，而是要把它作为一种产生于指挥关系中的证据。战果、兵数和伤亡若无多种材料约束，尤其不宜直接当作精确事实；一座城被报称攻下，也应与守军是否持续、市场能否恢复、周边道路谁在使用等问题分开。

清廷在内战中使用的八旗、绿营、勇营和地方团练，名目不同，招募、饷给、统属和驻扎的实际也不相同。不能因为都被称为“清军”就假定其纪律、装备和同居民的关系一致，更不能把勇营的亲缘、乡缘或师承网络写成无条件的忠诚。军队内部有上级与下属的责任关系，也有薪饷、升迁、伤病、家属和乡里牵出的要求。地方居民面对驻军时，可能从中谋得雇佣、保护或交易机会，也可能承受借宿、征购和暴力；同一支军队在不同地点的行为，需要具体文书和地方材料说明。国家军事能力正是在这些不稳定的关系中被维持，而不是在编制名称出现时已经完成。

太平天国方面的制度与动员亦须在自己的材料和行动范围内理解。天京的诏令、宗教语言、官制与军事部署说明一个政权试图如何组织权威，却不能只以清方所称的“贼”或后来的成败来解释其支持者和反对者。太平军控制区内的征粮、土地、性别规范、城市管理和军事迁徙，在不同时间与地点有不同的实践；对手的军报、租界观察与太平文书各有强烈立场。研究者可以比较这些材料对同一地点的称谓、日期和行动描述，却不应在缺乏地方记录时替居民补写明确的政治动机。有人加入、顺从、离开或与各方交易，可能是信念、亲属、暴力威胁、生计和局势判断交织的结果，不能被还原为单一选择。

上海在1853年小刀会起事及其后战事中的城市变化，也展示了华人与租界空间并非彼此隔绝。人口进入租界、商号调整经营、地方官与外国领事交涉、航运和雇佣改变，能够在不同类型的材料中留下痕迹；但“避入租界”不等于每个人得到同等安全，更不等于租界成为脱离周边战争的岛屿。房租、劳工、卫生、治安和粮食供应都使城市内部产生新的压力。外国报纸与领事档案对这些变化的记述往往集中于本国侨民、商业利益和治安担忧，中文地方材料则有另一套行政语言。把它们并读可看见城市空间如何被重新划分，却仍要避免把一座港市的经验概括为整个中国社会的转向。

战争与环境交织时，时间尺度尤其容易混乱。一次决口、一次旱灾或疫病流行可以造成紧迫的冲击，但堤防维护、土地淤积、人口流动和市场恢复往往持续多年。气候重建、水文研究和地方灾异记录能为这一长时段提供约束，却不能单独说明某户为何离乡、某处为何发生抢粮，或某军为何失利。反之，个别灾民的陈述虽有强烈的经验密度，也不能代表区域降雨或粮食总量。将环境材料与行政、市场和社会材料放在同一段叙述中，关键不是寻找一个万能原因，而是明确它们各自能说明的尺度：自然现象、行政认定、救济动作与人群后果必须逐项区分。

1870年代的义赈网络值得从这种区分中理解。士绅、商人、官员、报人、传教士和海外捐助者所筹集的银米，构成跨地区关切与新式信息传播的一个面向；但募捐公告、善款总数和名人名单并不自动等于粮食送到最急需之处。筹款、购买、转运、地方接收、名册登记和实际分配之间，每一环都会受价格、交通、地方权力和信息不均的影响。义赈材料能够让人追踪组织者如何设想灾区、如何证明自身可信，却也常使未被登记者难以出现。将义赈单纯赞为民间救国，或单纯斥为精英表演，都不如追问它在何处形成了可验证的救助，何处又留下了无法弥补的缺口。

西北与新疆的道路还提醒我们，跨境贸易并不在战争开始时消失、在条约签订后才重新出现。商队、翻译、向导、牲畜商和地方经纪人会依据关卡、税费、战事传闻和邻近市场调整路线；他们的活动把不同政治实体中的价格与消息联系起来，却不使边界失去意义。俄国、中亚和英印方面的观察材料常对贸易与军事动向格外敏感，也会把自己的战略焦虑投射到地方。清朝的档案对关防、贡道和军台的关切则不必等于绿洲居民或牧民的优先次序。跨语言地名和人名的对读可以减少误认，但不应强行制造一一对应；当史料无法确认某一地点或日期时，保留不确定性比用后来的地图填补更合乎证据。

边地的水利是另一种不宜被军事叙述吞没的基础设施。绿洲的渠水、堰坝、田地与劳役安排关系到城镇能否供给人口和驻军，也关系到地方权威如何被日常实践。军队占领或设官可以改变征税和劳役的要求，但并不等于能够立即掌握水源的分配、维修知识和社区协商。有关水利的地方文书、测绘和考古遗存保存有限，且常难以直接连接某次战役；正因如此，不能用一份凯旋叙事代替长期生计。讨论“经营新疆”时，至少应把军事推进、行政设置、移民安置、水利维护和贸易恢复区分为不同问题，而不是把它们折叠为一个胜利动词。

1877年前后的财政压力也显示，帝国并没有一个可随时从容支配的统一钱袋。海关收入、盐课、地丁、厘金、捐输、借款和地方自筹各有征收范围、时间和信用条件。用于西征、河工、赈济或厂局的款项在文书上可以被指定，却仍可能遭逢折色、拖欠、挪移和路径成本。财政史最有意义之处，不在于将每一笔钱化为抽象的国家能力，而在于揭示资源从谁处取得、经谁转手、在何处停留，以及风险由谁承担。数字与预算若没有地域、单位和会计口径，只能作为模糊的量感，不能支撑精确比较；同样，某项支出的高低也不能单独解释一场战争或一项改革的结果。

因此，1851至1877年的历史既不能由“中央衰弱、地方坐大”一句收束，也不能由“危机促进近代化”一句收束。前一种说法忽略中央文书、任命和财政框架仍在塑造地方行动，后一种说法则容易把战争的破坏、强制与不平等转换成进步的前提。更接近材料的理解，是看到各种权力在持续协商、竞争和失败中重新组合：有的官员和将领获得更大的调度空间，有的商人和中介进入新的契约关系，有的地方组织承担救济与防卫，有的家庭失去土地、劳力和发言渠道。结构变化确实发生，但它们从不以同一速度或同一代价抵达所有地方。

\section{从记录到判断}

在这样一段多线并进的历史里，编年仍然重要，但编年不是把所有材料压成单一时间表。1851年的广西起事、1853年的南京易手、1855年的黄河决口、1856年以后英法战争、1864年天京陷落、1860年代西北与云南的冲突、1871年伊犁危机及1874年的台湾交涉，都可以由相对可靠的近时记录定位；它们之间的关系却必须由道路、财政、人员与地方材料逐层检验。一件事在年表上相邻，不等于一方的命令立刻造成另一方的结果。反过来，相隔数年的政策也可能因债务、迁移、河道或机构延续而在后来留下后果。叙事的任务是说明可追索的联结，同时标出没有证据能跨越的空隙。

材料的沉默并不只存在于基层。朝廷档案可能缺失附件、抄录或地方回报，地方志可能在战乱中散佚或在重修时重排记忆，外文材料也可能因翻译、选编和档案分类而失去原来的语境。不同材料彼此矛盾并不必然说明其中一方无用，往往正显示它们服务于不同的行政、商业、军事或纪念目的。对日期、地点和行动的核对，应尽量寻找独立来源；对动机与感受，则更应节制，用材料真正允许的范围说话。没有居民自述时，不能把官员的安抚语言当作居民的认同；没有交接记录时，不能把预算当作到款；没有长期观察时，不能把一次攻城当作稳定秩序。

这种证据边界也使比较成为可能。比较长江与西北，不是问哪一处更“重要”，而是看水运与驮运、商路与军台、口岸与绿洲如何分别限定征发和信息。比较天津与台湾，也不是把两处都归入外交史便停止，而是辨认条约文本、领事交涉、地方防务和社群接触各在何处发生。比较战后江南与灾荒华北，则可看出复垦、赈济和河工为何需要不同材料尺度。只有承认差异，才不会把帝国写成一台无处不在的机器，或把地方写成同中央毫无关联的孤岛。

最后，战争留下的并不只是可供后人计算的制度遗产。它也留下被毁的住处、延误的婚期、未能归还的债务、改变的职业、失散的亲属和难以写入档案的恐惧。史料不能完整复原这些经验，写作也不应以想象填满它们；但正因为知道记录的限度，才应避免让胜利、条约、设局或复垦的宏大词语替它们结案。将政策与执行、战场与家庭、边界文本与地方道路分开考察，是为了在不夸大证据的前提下，保留那些被长时段危机改变的生活。

\noindent 这种谨慎并不要求叙事退回到彼此无关的片段。恰恰相反，只有把征发、运输、迁徙、交涉和重建放在可核查的关系中，才能看见一地的选择如何改变另一地的负担。历史中的联系不是预先铺好的轨道，而是在船只、文书、银两、市场和人的行动中暂时形成；它既可能把资源送到前线，也可能把风险转给无法拒绝的人。对这些联系的描述应当具体到能够被材料支持的程度，并为仍然不可知的部分留下位置。
\begin{sourcenote}[同一场危机为何会留下彼此不相同的地图]
《清文宗实录》《清穆宗实录》、军机处档、战事方略与地方奏报，最适合追索朝廷何时获报、命令怎样发出以及奏报者希望上级看见什么；它们把道路写成军需和行政问题，常不记录一处市场被征用后如何继续生活。《筹办夷务始末》、中外照会和条约文本能钉住谈判的文句、程序与一方提出的权利，却不能因同属官方或外交文书便彼此印证为执行结果。太平天国文书、地方档案、日记、租界报刊和外国军事报告改变了观察的站位和语言，而非自动补齐沉默：并读它们可以缩小对路线、消息和行动的误判，仍不能单独统计死伤、证明每一地的控制，或概括参战者与居民的意图。
\end{sourcenote}

\noindent 以下两组叙事从这些相互牵动而不同步的节点进入：先追问长江、黄河和省际筹饷怎样让战争持续，再回到西北、天津和台湾，观察重建、边路与条约网络如何把同一帝国拉向不同方向。

\section{多场战争挤在同一条粮道上}

1851年金田起事后，太平军北上，1853年占领南京并建都天京。长江城市、粮道和税源随之成为战场；北伐、西征、上海小刀会、捻军扩展又把华北、江淮和租界城市带进不同的战争节奏。1856年天京事变削弱太平天国的协调能力，同年第二次鸦片战争开始，清廷不得不在内战与海防之间分配本已不足的兵饷和注意力。1858年条约、1860年北京危机不是内战的插曲，而是重新规定财政、港口和外交机构的压力来源。

危机的关键不只是“中央弱、地方强”。湘军、淮军、团练与地方筹饷是在训练、家族网络、绅商捐输、厘金和军需转运中形成的；它们提高了清廷在若干战区的作战能力，也让军队、税收和任命更深地嵌入省级关系。1864年天京陷落终结了太平天国的中心政权，却没有使难民、荒地、债务和地方驻军同时消失。1868年捻军主力被击败后，战争留下的筹饷制度仍在塑造国家。

1855年黄河改道尤能显示战争并非纯粹的军事故事。河道、堤防、迁徙、粮价与军队移动相互影响；铜瓦厢的一次决口既不能解释全部战局，也不能被放在战事之外。地方家庭要在征发、逃难、守土、借贷和重建之间作出有限选择，正是这些选择使省份成为战后政治的基础。

\subsection{材料与争议}

军机档、方略和实录可追踪任命、城池和清廷的军费命令；太平文献、地方志、契约、日记、租界报刊和传教士材料各有立场与地域限制。兵数、伤亡、地方支持和战争损失不宜接受任何一方战报的总计，应将军事行动与社会后果分层叙述。

\section{重建不是回到原处}

战后的省份首先面对的不是一个抽象的“恢复”，而是河道、仓储、道路和城镇能否再次把人、粮和消息接在一起。地方官要呈报灾情和军需，商人、船户与地方组织要判断哪一条路线仍可通行，离散者则在征发、借贷、避难和返乡之间寻找余地。朝廷能够发出命令，并不表示银粮已经抵达，或一座城周围的市场已经重新运转。省级军政网络正是在这些断裂处获得作用，也把筹饷、运输和治安的成本留在地方社会。

这一章所能追踪的，是几个危机节点如何改变了重建的条件，而不是为所有地区写出同一条复原曲线。实录、军机档和督抚呈报较能保存朝廷获知何事、怎样下令以及哪些资源被要求调动；它们对远处道路是否畅通、征购怎样落到家庭、谁能离开或留下，常常说得很少。地方志、账簿、契约、日记和报刊可以补出一些地点的断面，却又受保存、书写者和城市视野限制。于是，下文把文书中的行政语言同河流、城镇和边路的实际约束分开，也不把任何一方报告中的规模、损失或控制范围当作不待检验的结论。

\subsection{长江失去一个枢纽后，省份怎样筹措}

\eventanchor{EQ-1853-NANJING-TAKING}{太平军攻占南京并改名天京（1853）}，使长江上一座大城成为长期政权中心。城名的改变并不只是一项象征：江面通行、周边府县的粮食往来、商路的风险和守城所需的物资，都要随之重新安排。对清廷而言，原先可以借助的江南财政与文化枢纽已成为必须争夺和绕开的空间；对沿江居民而言，谁在何处征发、谁能保护船只和货物，往往比远处的政权名称更直接。

朝廷的战事记录和地方呈报能确认南京易手及其后的军事压力，却不能替每一段江路说明货物是否照常流动，或每一处乡村如何筹措军粮。把“长江税源”写成一个可以随城池一并转移的整体，反而会抹去口岸、渡口、乡镇和省界之间的差异。战时筹饷之所以越来越依赖地方中介，不是因为某一道命令突然废除了中央财政，而是因为征解、转运和验收都要穿过不稳定的水陆路径。

\subsection{一条新河道与一个移动战场}

\eventanchor{EQ-1855-YELLOW-RIVER-COURSE-SHIFT}{铜瓦厢决口、黄河改道（1855）}把华北的堤防、耕地、聚落和运输条件推入新的不确定性。决口及河道北流是可以定位的节点；至于不同河段何时受淹、迁徙持续多久、粮价怎样变化，则不能从一份奏报推及整条黄河。河工文书记录的是修防、请款和呈报的行政过程，不等于堤防即时奏效；地方材料即使留下灾情，也往往只照见保存下来的县份和书写者。

同年，\eventanchor{EQ-1855-NIAN-REBELLION-EXPANSION}{捻军活动扩大（1855）}，又使淮北、安徽和华北的村落、堤防与商道进入另一种战争节奏。灾荒、流民、盐业与地方武装网络为行动提供了条件，但不能把水患直接写成一支武装的单一原因。骑兵、团练、运粮者和避难者各自沿着可走的路寻找资源或安全；他们的路线有时相交，有时彼此阻断。官府以“匪”概括的对象并不因此成为一个目的相同的群体，地方材料所能保留的，也只是这些选择的零散痕迹。

两件发生在同一年的事没有一条自动相连的因果线，却共同压缩了财政和运输的余地。河道变化使既有水路与堤工需要重新估量，移动战场又使米粮、牲畜和劳力更难按原有路径调度。省份此后不是被动承受中央压力的容器：地方官、绅商、团练与居民在可用道路和有限信用之间作出的安排，构成了清廷此后能够继续行动的条件，同时也使负担的去向更不均衡。

\subsection{天京的裂口与地方网络的代价}

\eventanchor{EQ-1856-TIANJING-INCIDENT}{天京事变（1856）}在太平天国核心内部造成权力冲突和杀戮，削弱了其协调军政资源的能力。现存的清方记录、奏折与后出的叙述能够提示这一裂口的存在及其政治后果，却不能为城内每一项行动、每一位死者或每一种动机给出无争议的数目和解释。将它写作一场立即决定所有战局的宫廷戏剧，同样会遮住沿江军队、乡村和对手在消息迟缓、补给不足中作出的不同判断。

裂口之所以与省份重建有关，在于军队和行政并不会因首都内部的冲突而停止索取粮饷。谁能筹款、募勇、保住一段航道或维持一座城周边的交换，便在战局中取得新的位置；但这种能力也以地方借贷、捐输、征购和长期驻军为代价。后来的省级军政网络因而既是重建的工具，也是把战争遗留的财政义务和武装关系嵌入地方的渠道。它解决的是协调资源的迫切问题，不能据此推断所有地区已经恢复秩序。

\subsection{京师危机所暴露的运输上限}

\eventanchor{EQ-1860-BEIJING-CAMPAIGN}{英法联军进逼北京（1860）}时，天津一线的谈判、京师防务和热河方向的行幸在短期内相互挤压。圆明园遭焚掠与咸丰帝出走是清楚的节点；英法、清廷和地方观察者留下的记录却服务于不同的军事、外交和政治目的。它们能够说明危机怎样被陈述和处理，不能仅凭一方的叙事便量定各处兵力、开支、掠夺范围或地方反应。

北京的危机也不应被写成压倒内地一切经验的单一中心事件。它与黄河改道、长江战场之间没有共同的即时因果，却都迫使清廷在有限的船路、驿递、饷项和官员注意力中作选择。海路和天津的交涉增加了信息与外交的压力，内战地区则继续要求粮秣和转运。能够抵达北京的呈报，未必等于同一消息已经抵达沿线州县；能够谈定的文本，也不等于港口、道路或省级财政已按其条款改变。重建于是始终受制于运输与信息的速度。

\subsection{西北不是内地战局的余波}

\eventanchor{EQ-1865-DUNGAN-REVOLT}{陕西、甘肃等地战事扩展（1865）}后，西北的城镇、商路、宗教社群和地方军队被卷入多年冲突。地方冲突、军政失序和官府的族群分类相互激化，是理解战争条件的起点，却不足以把不同地方的行动者归入一个单一阵营。官书中的“回”等称谓首先属于行政和征剿的语言；它不能替代人们的自称、宗教联系、市场位置或在战争中的具体选择。

这里的边地不是等候中央重新填补的空白。陕西、甘肃与通往更西方的道路，连接着城镇、绿洲、牧地和跨境贸易的多种节奏；一段道路的失守或重开，会分别影响行军、粮运、逃亡和商业，却未必在所有地点产生同样结果。清廷的军机文书和战事汇编能够追踪调兵、命令和清方所声称的进展；要判断社区边界如何变化、市场如何中断、迁徙和暴力留下何种后果，仍须并读地方材料以及满、蒙、维吾尔、俄文和其他能够取得的区域材料。许多这样的材料保存不全、译名未必相合，正是我们不能把西北战争压缩成一张中央战报的原因。

从省份财政看，西北战争再次显出距离的成本。饷项不是一个从京师发出便自然抵达前线的数目，而是一连串筹措、折算、押运、采购和交接；任何一个环节受阻，前线和沿线社会承受的压力都会改变。现有材料足以说明长期战争要求跨省调度，不能据此替每一处征发、每一批粮饷或每一条边路断言执行的程度。所谓重建，在这里首先是重新取得有限的通行和供给能力，而非对地方生活的即时掌握。

\subsection{天津：地方冲突如何进入条约网络}

\eventanchor{EQ-1870-TIANJIN-MASSACRE}{天津教案（1870）}让育婴堂、传教机构、地方传闻、官府处置和法国外交照会在同一座城市相互碰撞。法国领事、修女和中国民众的死伤使冲突迅速进入外交交涉，但不能因此把天津居民、教民、传教士或官员各自写成没有分歧的阵营。谁相信何种传闻、谁试图调停、谁在暴力中失去行动空间，须回到文本产生的地点和目的中辨认。

\begin{event}{天津教案：地方传闻何以变成外交危机}{同治九年（一八七〇年）}{天津民众、天主教机构、清廷与法国}{冲突及外交后果可核；死伤、谣言传播和个别行动者动机在中外材料中差异很大}
总理衙门档、地方呈报与法国外交材料可以互校危机如何被报告和追责。照会所能确证的是国家怎样提出责任与要求；地方档案和教会通信则分别保留城市社区与宗教网络的片段。它们对人数、责任与意图并不一致，也都不能自行代替那些没有留下文字的人。

天津使战后重建面对另一种限制：省级官员处理地方治安时，已经必须估量条约关系、领事保护和跨国信息的后果。可是，外交文本不能反过来证明城市秩序已经被外来制度接管。保留这种不一致，才能看见教案既是地方社会的撕裂，也是交通、报信和条约网络如何改变行政选择的事件。
\end{event}

\subsection{制度得失：能调动，不等于能复原}

从南京到天津、从黄河到西北，清廷重新取得行动能力的方式往往是借助省级筹饷、地方军政和跨区域运输。这样的安排能在危机中把命令、物资和人手连接起来，却把借贷、征购、驻军和道路风险分配给不同地区与人群。它并未把帝国恢复为战前的样子，而是形成了一套必须依靠地方节点、港口信息和边路通行才能运作的秩序。

这也是证据最容易诱发过度判断之处。制度章程、战报或外交照会能让我们看见计划、要求和国家间的措辞；地方执行、实际产出和普通人的经验需要另一类材料，且常常无法完整取得。读者应把省份的“重建”理解为持续的、地域不均的协商过程：有些路被重新接上，有些社区仍在战争和迁徙的余波中，而边地的行动者从未只是中央叙事的背景。

